% ============================================================
%  PEMBAHASAN  --  MATEMATIKA  --  SM ITS -- Paket 7.2
%  45 Soal (Set Latihan 3 + Set Latihan 2)
% ============================================================

\documentclass[11pt]{article}
\usepackage[utf8]{inputenc}
\usepackage[T1]{fontenc}
\usepackage[a4paper,margin=2.5cm,top=2.8cm]{geometry}
\usepackage{amsmath,amssymb}
\usepackage{graphicx}
\usepackage{enumitem}
\usepackage{multicol}
\usepackage{xcolor}
\usepackage[most]{tcolorbox}
\usepackage{fancyhdr}
\usepackage{booktabs}
\usepackage{icomma}
\usepackage{array}
\usepackage{tikz}

\definecolor{mapelcolor}{RGB}{20,110,60}
\definecolor{mapelbg}{RGB}{228,246,234}
\definecolor{pembcolor}{RGB}{100,50,130}
\definecolor{pembbg}{RGB}{245,238,252}
\definecolor{gold}{RGB}{210,160,30}

\pagestyle{fancy}\fancyhf{}
\lhead{\small SM ITS -- Paket 7.2 -- Pembahasan}
\rhead{\small Matematika}
\cfoot{\small\thepage}
\renewcommand{\headrulewidth}{0.4pt}

\newtcolorbox{soal}[1]{
  breakable,
  colback=mapelbg, colframe=mapelcolor,
  coltitle=white, colbacktitle=mapelcolor,
  fonttitle=\bfseries\sffamily,
  title=Nomor #1,
  boxrule=0.8pt, arc=3pt,
  left=3mm, right=3mm, top=2mm, bottom=2mm}

\newtcolorbox{pembox}{
  breakable,
  colback=pembbg, colframe=pembcolor,
  coltitle=white, colbacktitle=pembcolor,
  fonttitle=\bfseries\sffamily,
  title=Pembahasan,
  boxrule=0.8pt, arc=3pt,
  left=3mm, right=3mm, top=2mm, bottom=2mm}

\newtcolorbox{catatan}{
  breakable,
  colback=gold!15, colframe=gold!80!black,
  coltitle=white, colbacktitle=gold!80!black,
  fonttitle=\bfseries\sffamily,
  title=Catatan,
  boxrule=0.8pt, arc=3pt,
  left=3mm, right=3mm, top=2mm, bottom=2mm}

\newcommand{\jawaban}[1]{\textbf{Jawaban:}\enspace\colorbox{gold!30}{\textbf{#1}}}
\newcommand{\konsep}[1]{\textit{Konsep: #1.}}
\newenvironment{pil}{%
  \begin{enumerate}[label=(\alph*),leftmargin=*,itemsep=1pt,topsep=3pt]}%
  {\end{enumerate}}
\linespread{1.05}

\begin{document}
\thispagestyle{empty}
\begin{center}
  {\LARGE\bfseries PEMBAHASAN MATEMATIKA}\\[0.4em]
  {\large SM ITS -- Paket 7.2}
\end{center}

\vspace{0.4cm}\noindent\textbf{Kunci Jawaban}

\vspace{0.3em}{\small
\begin{tabular}{@{}*{10}{c@{\hspace{1.1em}}}@{}}
\textbf{1}(c)&\textbf{2}(c)&\textbf{3}(b)&\textbf{4}(a)&\textbf{5}(c)&
\textbf{6}(e)&\textbf{7}(b)&\textbf{8}(d)&\textbf{9}(e)&\textbf{10}(e)\\
\end{tabular}

\vspace{0.2em}
\begin{tabular}{@{}*{10}{c@{\hspace{1.1em}}}@{}}
\textbf{11}(e)&\textbf{12}(c)&\textbf{13}(b)&\textbf{14}(d)&\textbf{15}(c)&
\textbf{16}(c)&\textbf{17}(d)&\textbf{18}(c)&\textbf{19}(a)&\textbf{20}(e)\\
\end{tabular}

\vspace{0.2em}
\begin{tabular}{@{}*{10}{c@{\hspace{1.1em}}}@{}}
\textbf{21}(d)&\textbf{22}(d)&\textbf{23}(d)&\textbf{24}(d)&\textbf{25}(e)&
\textbf{26}(b)&\textbf{27}(c)&\textbf{28}(b)&\textbf{29}(e)&\textbf{30}(d)\\
\end{tabular}

\vspace{0.2em}
\begin{tabular}{@{}*{10}{c@{\hspace{1.1em}}}@{}}
\textbf{31}(c)&\textbf{32}(d)&\textbf{33}(c)&\textbf{34}(b)&\textbf{35}(a)&
\textbf{36}(a)&\textbf{37}(a)&\textbf{38}(d)&\textbf{39}(d)&\textbf{40}(d)\\
\end{tabular}

\vspace{0.2em}
\begin{tabular}{@{}*{5}{c@{\hspace{1.1em}}}@{}}
\textbf{41}(a)&\textbf{42}(a)&\textbf{43}(a)&\textbf{44}(b)&\textbf{45}(b)\\
\end{tabular}
}

\clearpage

% ===================================================================
%  BAGIAN A: SET LATIHAN 3  (Q1--25)
% ===================================================================
{\noindent\large\bfseries Bagian A \textendash\ Set Latihan 3\par}\smallskip\hrule\medskip

\begin{soal}{1 -- Integral Fungsi Kasus}
$\int_{-1}^1 f(x)\,dx$ untuk fungsi kasus.
\end{soal}
\begin{pembox}
$\int_{-1}^0(x+1)\,dx+\int_0^1\cos\pi x\,dx$.
Bagian 1: $\left[\frac{x^2}{2}+x\right]_{-1}^0=0-(\frac12-1)=\frac12$.
Bagian 2: $\left[\frac{\sin\pi x}{\pi}\right]_0^1=0$.
Total $=\frac12$.
\jawaban{(c) $\dfrac12$}
\end{pembox}

\begin{soal}{2 -- Turunan Trigonometri}
$f(x)=x\sin(2x^2)$. Nilai $f'\!\left(\sqrt{\pi/12}\right)$.
\end{soal}
\begin{pembox}
Aturan hasil kali dan rantai: $f'(x)=\sin(2x^2)+x\cdot\cos(2x^2)\cdot4x=\sin(2x^2)+4x^2\cos(2x^2)$.

Untuk $x=\sqrt{\pi/12}$: $x^2=\pi/12$, $2x^2=\pi/6$.
$f'=\sin\frac\pi6+4\cdot\frac\pi{12}\cos\frac\pi6=\frac12+\frac\pi3\cdot\frac{\sqrt3}{2}=\frac12+\frac{\sqrt3\pi}{6}=\dfrac{3+\sqrt3\pi}{6}$.
\jawaban{(c) $\dfrac{3+\sqrt3\pi}{6}$}
\end{pembox}

\begin{soal}{3 -- Kontinuitas}
$f$ kontinu di $x=0$; cari $k$.
\end{soal}
\begin{pembox}
Dari kiri: $f(0)=1+k$. Dari kanan: $\lim_{x\to0^+}\frac{3\sin2x}{x}=3\cdot2=6$.
$1+k=6\Rightarrow k=5$.
\jawaban{(b) $5$}
\end{pembox}

\begin{soal}{4 -- Bentuk Akar}
$\sqrt{x/3}+\sqrt{x/27}+\sqrt{x/243}=\frac{156}{54}$.
\end{soal}
\begin{pembox}
Faktorkan $\sqrt x$: $\sqrt x\cdot\frac{1+\frac13+\frac19}{\sqrt3}=\frac{156}{54}$.

Isi kurung $=\frac{13}{9}$; $\frac{13}{9\sqrt3}\sqrt x=\frac{26}{9}$. $\sqrt x=2\sqrt3\Rightarrow x=12$.
\jawaban{(a) $12$}
\end{pembox}

\begin{soal}{5 -- Logaritma}
$\log_x81=8$, $\log_3x=y$. Nilai $xy$.
\end{soal}
\begin{pembox}
$x^8=81=3^4\Rightarrow x=3^{1/2}$. $y=\log_3(3^{1/2})=\frac12$. $xy=\sqrt3\cdot\frac12=\frac{\sqrt3}{2}$.
\jawaban{(c) $\dfrac{\sqrt3}{2}$}
\end{pembox}

\begin{soal}{6 -- Garis Singgung}
$y=x^4-x^3-2x^2+5x+1$ di $(1,4)$. Nilai $a+b$.
\end{soal}
\begin{pembox}
$y'=4x^3-3x^2-4x+5$; $y'(1)=4-3-4+5=2$, sehingga $a=2$.
$b=4-2(1)=2$. $a+b=4$.
\jawaban{(e) $4$}
\end{pembox}

\begin{soal}{7 -- Vektor dan Sudut}
$|\vec a+\vec u|^2=27$, $|\vec a|=5$, $|\vec u|=1$.
\end{soal}
\begin{pembox}
$|\vec a+\vec u|^2=25+1+2\vec a\cdot\vec u=27\Rightarrow\vec a\cdot\vec u=\frac12$.
$\cos\theta=\frac{1/2}{5}=\frac{1}{10}$; $\tan\theta=\sqrt{1-\frac{1}{100}}\big/\frac{1}{10}=\sqrt{99}$.
\jawaban{(b) $\sqrt{99}$}
\end{pembox}

\begin{soal}{8 -- Lingkaran Dalam Persegi}
$P(-5,0)$, $Q(1,0)$, sisi $PQ=6$, lingkaran di tengah persegi.
\end{soal}
\begin{pembox}
Pusat lingkaran: tengah persegi $=\left(\frac{-5+1}{2},\frac{0+6}{2}\right)=(-2,3)$; $r=3$.
$(x+2)^2+(y-3)^2=9\Rightarrow x^2+y^2+4x-6y+4=0$.
\jawaban{(d) $x^2+y^2+4x-6y+4=0$}
\end{pembox}

\begin{soal}{9 -- Statistika}
13 bilangan, rata-rata 3; dua dikeluarkan, diganti 7; rata-rata baru 4.
\end{soal}
\begin{pembox}
Jumlah semula $=39$. Setelah penggantian: $39-S+7=48\Rightarrow S=-2$.
Pasangan dengan jumlah $-2$: $-1{,}5$ dan $-0{,}5$.
\jawaban{(e) $-1{,}5$ dan $-0{,}5$}
\end{pembox}

\begin{soal}{10 -- Barisan Aritmetika}
$a_1+a_5+a_{10}+a_{15}+a_{20}+a_{24}=225$. Jumlah $S_{24}$.
\end{soal}
\begin{pembox}
Jumlah $=6a_1+69d=225$. $S_{24}=12(2a_1+23d)=4(6a_1+69d)=4\times225=900$.
\jawaban{(e) $900$}
\end{pembox}

\begin{soal}{11 -- Kecepatan Relatif}
James menyusul Karen setelah $t$ menit.
\end{soal}
\begin{pembox}
$vt=0{,}8v(t+15)\Rightarrow0{,}2t=12\Rightarrow t=60$.
\jawaban{(e) $60$ menit}
\end{pembox}

\begin{soal}{12 -- Sistem Tak Punya Solusi}
Hasil kali nilai $b$.
\end{soal}
\begin{pembox}
$\det\begin{pmatrix}2&5&1\\-4&b&6\\0&-3&-b\end{pmatrix}=-2(b-2)(b+12)=0\Rightarrow b=2$ atau $b=-12$.

Hasil kali $=2\times(-12)=-24$.
\jawaban{(c) $-24$}
\end{pembox}

\begin{soal}{13 -- Vektor Tegak Lurus}
$|\vec a+\vec b|=|\vec a|$. Pernyataan yang benar.
\end{soal}
\begin{pembox}
Kuadratkan: $|\vec a|^2+2\vec a\cdot\vec b+|\vec b|^2=|\vec a|^2\Rightarrow2\vec a\cdot\vec b+\vec b\cdot\vec b=0$.
$(2\vec a+\vec b)\cdot\vec b=2\vec a\cdot\vec b+\vec b\cdot\vec b=0\Rightarrow2\vec a+\vec b\perp\vec b$.
\jawaban{(b) $2\vec a+\vec b$ dan $\vec b$ saling tegak lurus}
\end{pembox}

\begin{soal}{14 -- Komposisi Fungsi}
$f(0)+g(0)$.
\end{soal}
\begin{pembox}
$f(t)=3at^2+b$; $f(\sqrt3)=9a+b=38$. $g(t)=\frac{b}{5}t^2+a$; $g(5\sqrt2)=10b+a=24$.
Dari sistem: $a=4$, $b=2$. $f(0)=b=2$, $g(0)=a=4$. $f(0)+g(0)=6$.
\jawaban{(d) $6$}
\end{pembox}

\begin{soal}{15 -- Deret Tak Hingga}
$\sum(2^n+5^n)/10^n$.
\end{soal}
\begin{pembox}
$\sum(\frac12)^n+\sum(\frac15)^n=\frac{1/2}{1-1/2}+\frac{1/5}{1-1/5}=1+\frac14=\frac54$.
\jawaban{(c) $\dfrac54$}
\end{pembox}

\begin{soal}{16 -- Peluang}
Anto di tim berisi 3 dari 6; peluang Mika bersama Anto.
\end{soal}
\begin{pembox}
Tim Anto butuh 2 dari 5 orang lain; peluang Mika terpilih $=\frac25=40\%$.
\jawaban{(c) $40\%$}
\end{pembox}

\begin{soal}{17 -- Invers Komposisi}
$(f\circ g^{-1})^{-1}(-2)$.
\end{soal}
\begin{pembox}
$f(t)=\frac{3t-1}{2}\Rightarrow f^{-1}(u)=\frac{2u+1}{3}$. $g^{-1}(y)=\frac{5y+66}{4}\Rightarrow g(z)=\frac{4z-66}{5}$.

$(f\circ g^{-1})^{-1}=g\circ f^{-1}$. $f^{-1}(-2)=\frac{-3}{3}=-1$. $g(-1)=\frac{-4-66}{5}=-14$.
\jawaban{(d) $-14$}
\end{pembox}

\begin{soal}{18 -- Suku Banyak}
$2x^3-px^2+4x+q$ habis dibagi $2x^2+x-1$. Jumlah akar.
\end{soal}
\begin{pembox}
$(2x-1)(x+1)$ $\Rightarrow$ $x=\frac12$ dan $x=-1$ adalah akar.
$P(\frac12)=0$: $q=\frac p4-\frac94$. $P(-1)=0$: $q=p+6$. Selesaikan: $p=-11$.
Jumlah 3 akar $=\frac{p}{2}=-\frac{11}{2}$.
\jawaban{(c) $-\dfrac{11}{2}$}
\end{pembox}

\begin{soal}{19 -- Aturan Rantai}
$g(x)=[f(2x+1)+2]^2$. Nilai $g'(0)$.
\end{soal}
\begin{pembox}
$g'(x)=4[f(2x+1)+2]f'(2x+1)$.
$g'(0)=4[f(1)+2]f'(1)=4[-1+2](1)=4$.
\jawaban{(a) $4$}
\end{pembox}

\begin{soal}{20 -- Determinan}
$|2B|$.
\end{soal}
\begin{pembox}
$B=A\begin{pmatrix}3&0\\-1&4\end{pmatrix}$; $|B|=|A|\cdot12=60$. $|2B|=4|B|=240$.
\jawaban{(e) $240$}
\end{pembox}

\begin{soal}{21 -- Integral Definit}
$\int_{-2}^0(f(x)+2)\,dx$.
\end{soal}
\begin{pembox}
$\int_{-2}^0f=-p$. $\int_{-2}^0(f+2)=-p+4=4-p$.
\jawaban{(d) $4-p$}
\end{pembox}

\begin{soal}{22 -- Trigonometri}
$\sin\theta-\cos\theta=\frac13$. Nilai $\frac{\sin\theta}{\cos\theta}+\frac{\cos\theta}{\sin\theta}$.
\end{soal}
\begin{pembox}
Kuadratkan: $1-2\sin\theta\cos\theta=\frac19\Rightarrow\sin\theta\cos\theta=\frac49$.
$\frac{\sin^2\theta+\cos^2\theta}{\sin\theta\cos\theta}=\frac{1}{4/9}=\frac94$.
\jawaban{(d) $\dfrac94$}
\end{pembox}

\begin{soal}{23 -- Persamaan Kuadrat Baru}
Akar baru $x_1+1$ dan $x_2+1$ dari $2x^2-x-5=0$.
\end{soal}
\begin{pembox}
$x_1+x_2=\frac12$, $x_1x_2=-\frac52$. Akar baru: jumlah $=\frac52$, kali $=-\frac52+\frac12+1=-1$.
Persamaan: $2t^2-5t-2=0$.
\jawaban{(d) $2x^2-5x-2=0$}
\end{pembox}

\begin{soal}{24 -- Pertidaksamaan Nilai Mutlak}
$|{(x+3)}/{(x-4)}|<2$.
\end{soal}
\begin{pembox}
Kuadratkan: $(x+3)^2<4(x-4)^2\Rightarrow3x^2-38x+55>0\Rightarrow(3x-5)(x-11)>0$.
Solusi: $x<\frac53$ atau $x>11$.
\jawaban{(d) $x<\dfrac53$ atau $x>11$}
\end{pembox}

\begin{soal}{25 -- Pertidaksamaan}
$x^2<x$ dan $xy>y$. Yang selalu benar.
\end{soal}
\begin{pembox}
$x^2<x\Rightarrow0<x<1$. $xy>y\Rightarrow y(x-1)>0$; karena $x-1<0$, maka $y<0$.

Dengan $0<x<1$ dan $y<0$: $3x>0$ dan $-5y>0$, sehingga $3x-5y>0$ selalu.
\begin{catatan}
Opsi (e) pada naskah sumber tertulis $3x-5y<0$. Tanda yang benar adalah $>0$.
\end{catatan}
\jawaban{(e) $3x-5y>0$}
\end{pembox}

% ===================================================================
%  BAGIAN B: SET LATIHAN 2  (Q26--45)
% ===================================================================
\medskip{\noindent\large\bfseries Bagian B \textendash\ Set Latihan 2\par}\smallskip\hrule\medskip

\begin{soal}{26 -- Bentuk Akar}
$\dfrac{\sqrt{10}-\sqrt5}{\sqrt{10}+\sqrt5}=a+b\sqrt2$.
\end{soal}
\begin{pembox}
Rasionalkan: $\frac{(\sqrt{10}-\sqrt5)^2}{5}=\frac{15-10\sqrt2}{5}=3-2\sqrt2$.
$a=3$, $b=-2$. $a+b=1$.
\jawaban{(b) $1$}
\end{pembox}

\begin{soal}{27 -- Komposisi Fungsi Terbalik}
$f(x)=\sqrt{x+1}$; $(f\circ g)(x)=2\sqrt{x-1}$. Cari $g(x)$.
\end{soal}
\begin{pembox}
$\sqrt{g(x)+1}=2\sqrt{x-1}$. Kuadratkan: $g(x)+1=4(x-1)\Rightarrow g(x)=4x-5$.
\jawaban{(c) $4x-5$}
\end{pembox}

\begin{soal}{28 -- Determinan Matriks}
$AX=B+A^T$. Nilai $\det(X)$.
\end{soal}
\begin{pembox}
$\det(X)=\frac{\det(B+A^T)}{\det(A)}$. $A^T=\begin{pmatrix}3&0\\2&5\end{pmatrix}$; $B+A^T=\begin{pmatrix}0&-1\\-15&5\end{pmatrix}$.
$\det(B+A^T)=0-15=-15$. $\det(A)=15$. $\det(X)=-1$.
\jawaban{(b) $-1$}
\end{pembox}

\begin{soal}{29 -- Sistem Pertidaksamaan}
$3-x>4$ dan $\frac12x<5$.
\end{soal}
\begin{pembox}
$3-x>4\Rightarrow x<-1$. $\frac12x<5\Rightarrow x<10$. Irisan: $x<-1$.
\begin{catatan}
Pilihan (a)--(d) tidak memuat $x<-1$; opsi (e) merupakan koreksi berdasarkan hasil perhitungan.
\end{catatan}
\jawaban{(e) $x<-1$}
\end{pembox}

\begin{soal}{30 -- Invers Fungsi Logaritma}
Invers $y=\log_{1/2}x$.
\end{soal}
\begin{pembox}
Invers dari $y=\log_a x$ adalah $y=a^x$. Dengan $a=\frac12$: inversnya $y=\left(\frac12\right)^x$.
\jawaban{(d) $y=\left(\dfrac12\right)^x$}
\end{pembox}

\begin{soal}{31 -- Diskriminan dan Optimasi}
$m,n>0$; nilai terkecil $2m+3n$.
\end{soal}
\begin{pembox}
$D_1\geq0$: $m^2\geq4n^2\Rightarrow m\geq2n$.
$D_2\geq0$: $16n^2\geq8m\Rightarrow m\leq2n^2$.
$2n\leq m\leq2n^2\Rightarrow n\geq1$. Ambil $n=1$, $m=2$: $2(2)+3(1)=7$.
\jawaban{(c) $7$}
\end{pembox}

\begin{soal}{32 -- Cayley-Hamilton}
$A^2-xA+yI=O$. Nilai $x+y$.
\end{soal}
\begin{pembox}
$\mathrm{tr}(A)=9$, $\det(A)=14$. Cayley-Hamilton: $x=9$, $y=14$. $x+y=23$.
\jawaban{(d) $23$}
\end{pembox}

\begin{soal}{33 -- Logaritma dan Lingkaran}
$r=\log p^2$, $K=\log q^4$. Nilai $\log_p q$.
\end{soal}
\begin{pembox}
$K=2\pi r\Rightarrow4\log q=2\pi\cdot2\log p\Rightarrow\frac{\log q}{\log p}=\pi\Rightarrow\log_p q=\pi$.
\jawaban{(c) $\pi$}
\end{pembox}

\begin{soal}{34 -- Sistem Persamaan}
$x+y+z=18$, $x^2+y^2+z^2=756$, $x^2=yz$. Nilai $x$.
\end{soal}
\begin{pembox}
$y+z=18-x$. $(y+z)^2=y^2+z^2+2yz=(756-x^2)+2x^2=756+x^2$.
$(18-x)^2=756+x^2\Rightarrow324-36x=756\Rightarrow x=-12$.
\jawaban{(b) $-12$}
\end{pembox}

\begin{soal}{35 -- Parabola dan Garis Singgung}
Garis singgung $y=x^2$ sejajar dengan $g$ (menghubungkan $(-a,a^2)$ dan $(3a,9a^2)$).
\end{soal}
\begin{pembox}
Kemiringan $g=\frac{9a^2-a^2}{4a}=2a$. Garis singgung $y=x^2$ di $x=a$: $y-a^2=2a(x-a)\Rightarrow y=2ax-a^2$.
Memotong sumbu-$y$ di $y=-a^2$.
\jawaban{(a) $-a^2$}
\end{pembox}

\begin{soal}{36 -- Statistika -- Nilai Terbesar}
20 bilangan bulat positif berbeda, rata-rata 20. Bilangan terbesar.
\end{soal}
\begin{pembox}
Jumlah $=400$. 19 bilangan terkecil: $1+2+\cdots+19=190$. Terbesar $=400-190=210$.
\jawaban{(a) $210$}
\end{pembox}

\begin{soal}{37 -- Optimasi Volume}
$x^2+4xh=192$; volume $V=x^2h$.
\end{soal}
\begin{pembox}
$h=\frac{192-x^2}{4x}$; $V=48x-\frac{x^3}{4}$. $V'=48-\frac{3x^2}{4}=0\Rightarrow x=8$, $h=4$.
$V_{\max}=64\times4=256$ cm$^3$.
\jawaban{(a) $256$ cm$^3$}
\end{pembox}

\begin{soal}{38 -- Transformasi Geometri}
Translasi $(4,5)$ ke $(1,3)$; lalu cermin ke $(2,3)$.
\end{soal}
\begin{pembox}
$(x,y)=(1-4,3-5)=(-3,-2)$. Garis cermin = sumbu simetri ruas $(-3,-2)-(2,3)$.
Titik tengah $(-\frac12,\frac12)$; kemiringan ruas $=1$; kemiringan cermin $=-1$.
$y-\frac12=-(x+\frac12)\Rightarrow y=-x$.
\jawaban{(d) $y=-x$}
\end{pembox}

\begin{soal}{39 -- Barisan Geometri}
$a+b+c=-14$, $abc=216$, $c<a$. Nilai $c$.
\end{soal}
\begin{pembox}
$b^3=216\Rightarrow b=6$. $a+c=-20$; $ac=36$. Akar $t^2+20t+36=(t+2)(t+18)=0$: $t=-2$ atau $t=-18$.
Karena $c<a$, $c=-18$.
\jawaban{(d) $-18$}
\end{pembox}

\begin{soal}{40 -- Trigonometri}
$m=\sin^2x\cot^2x+\cos^2x\tan^2x$. Nilai $(2m+3)^2$.
\end{soal}
\begin{pembox}
$\sin^2x\cot^2x=\cos^2x$; $\cos^2x\tan^2x=\sin^2x$. $m=1$. $(2+3)^2=25$.
\jawaban{(d) $25$}
\end{pembox}

\begin{soal}{41 -- Persamaan Trigonometri}
$\cos(x+10°)=\sin(2x-10°)$, $0°\leq x\leq180°$.
\end{soal}
\begin{pembox}
$\sin(2x-10°)=\cos(100°-2x)$. Maka $x+10=100-2x\Rightarrow x=30°$, atau $x+10=-(100-2x)+360°\Rightarrow x=110°$.
\jawaban{(a) $\{30°,110°\}$}
\end{pembox}

\begin{soal}{42 -- Persamaan Trigonometri}
$\sin x-\tan x-2\cos x+2=0$, $0<x<\pi/2$.
\end{soal}
\begin{pembox}
Kalikan $\cos x$: $(\cos x-1)(\sin x-2\cos x)=0$. Pada interval, $\cos x\neq1$; $\sin x=2\cos x\Rightarrow\tan x=2$.
$\sin x=\frac{2}{\sqrt5}=\frac{2\sqrt5}{5}$.
\jawaban{(a) $\left\{\dfrac{2\sqrt5}{5}\right\}$}
\end{pembox}

\begin{soal}{43 -- Logaritma}
$3\log a-2\log b=4$; $4\log a+3\log b=11$. Nilai $b^2-a$.
\end{soal}
\begin{pembox}
$u=\log a$, $v=\log b$: $3u-2v=4$ dan $4u+3v=11$. Selesaikan: $u=2$, $v=1$.
$a=100$, $b=10$. $b^2-a=100-100=0$.
\jawaban{(a) $0$}
\end{pembox}

\begin{soal}{44 -- Eksponen}
$\dfrac{2^{1/2}+2^{1/3}}{2^{-1/2}+2^{-1/3}}=4^x$.
\end{soal}
\begin{pembox}
$2^{-1/2}+2^{-1/3}=\frac{a+b}{ab}$ dengan $a=2^{1/2}$, $b=2^{1/3}$.
Ruas kiri $=ab=2^{1/2+1/3}=2^{5/6}$. $4^x=2^{2x}=2^{5/6}\Rightarrow x=\frac{5}{12}$.
\jawaban{(b) $\dfrac{5}{12}$}
\end{pembox}

\begin{soal}{45 -- Suku Banyak -- Sisa}
$p(1)=2$, $p(-1)=0$; $R_q(x)=2x$. Sisa $p+q$ dibagi $x^2-1$.
\end{soal}
\begin{pembox}
$R_p(x)=ax+b$: $a+b=2$, $-a+b=0\Rightarrow a=1$, $b=1$; $R_p=x+1$.
Sisa $p+q$: $(x+1)+2x=3x+1$.
\jawaban{(b) $3x+1$}
\end{pembox}

\end{document}
